\documentclass[11pt]{article}
\usepackage[T1]{fontenc}
\usepackage[utf8]{inputenc}
\usepackage{lmodern}
\usepackage[letterpaper,margin=0.78in]{geometry}
\usepackage{microtype}
\usepackage{booktabs,longtable,tabularx,array,multirow}
\usepackage{amsmath,amssymb}
\usepackage{xcolor}
\usepackage{enumitem}
\usepackage{titlesec}
\usepackage{fancyhdr}
\usepackage{lastpage}
\usepackage{hyperref}
\usepackage[most]{tcolorbox}
\usepackage{parskip}
\usepackage{needspace}
\usepackage{ragged2e}
\usepackage{setspace}

\definecolor{Navy}{HTML}{17365D}
\definecolor{Blue}{HTML}{2F75B5}
\definecolor{Teal}{HTML}{1F6D6A}
\definecolor{Gold}{HTML}{B38B2D}
\definecolor{Gray}{HTML}{5E6B75}
\definecolor{LightBlue}{HTML}{EEF4FA}
\definecolor{LightGold}{HTML}{FBF6E8}
\definecolor{LightGray}{HTML}{F4F5F7}
\definecolor{LightTeal}{HTML}{EDF7F6}

\hypersetup{colorlinks=true,linkcolor=Navy,urlcolor=Blue,citecolor=Teal,pdfauthor={OpenAI},pdfsubject={Value investing study and implementation materials}}
\setlist[itemize]{leftmargin=1.5em,itemsep=2pt,topsep=3pt}
\setlist[enumerate]{leftmargin=1.7em,itemsep=2pt,topsep=3pt}
\renewcommand{\arraystretch}{1.16}
\setlength{\tabcolsep}{5pt}
\titleformat{\section}{\Large\bfseries\color{Navy}}{\thesection}{0.6em}{}[\vspace{-0.25em}\color{Blue}\rule{\linewidth}{0.5pt}]
\titleformat{\subsection}{\large\bfseries\color{Teal}}{\thesubsection}{0.5em}{}
\titleformat{\subsubsection}{\normalsize\bfseries\color{Gray}}{\thesubsubsection}{0.5em}{}
\pagestyle{fancy}
\fancyhf{}
\fancyhead[L]{\small\color{Gray}\DocShortTitle}
\fancyhead[R]{\small\color{Gray}Value Investing Reference Set}
\fancyfoot[C]{\small\color{Gray}Page \thepage\ of \pageref*{LastPage}}
\renewcommand{\headrulewidth}{0.4pt}
\renewcommand{\footrulewidth}{0pt}

\newtcolorbox{principlebox}[1][]{enhanced,breakable,colback=LightBlue,colframe=Blue,boxrule=0.6pt,arc=2mm,left=2mm,right=2mm,top=1.5mm,bottom=1.5mm,title={#1},fonttitle=\bfseries}
\newtcolorbox{cautionbox}[1][]{enhanced,breakable,colback=LightGold,colframe=Gold,boxrule=0.6pt,arc=2mm,left=2mm,right=2mm,top=1.5mm,bottom=1.5mm,title={#1},fonttitle=\bfseries}
\newtcolorbox{modernbox}[1][]{enhanced,breakable,colback=LightTeal,colframe=Teal,boxrule=0.6pt,arc=2mm,left=2mm,right=2mm,top=1.5mm,bottom=1.5mm,title={#1},fonttitle=\bfseries}
\newtcolorbox{worksheetbox}[1][]{enhanced,breakable,colback=white,colframe=Gray,boxrule=0.45pt,arc=1mm,left=2mm,right=2mm,top=1.5mm,bottom=1.5mm,title={#1},fonttitle=\bfseries}
\newcommand{\checkbox}{\(\square\)}
\newcommand{\blankline}{\rule{0.96\linewidth}{0.35pt}}
\newcommand{\sourceRef}[1]{\textit{Source anchor: #1}}
\newcolumntype{Y}{>{\RaggedRight\arraybackslash}X}
\newcolumntype{C}[1]{>{\centering\arraybackslash}p{#1}}

\newcommand{\makedoccover}[4]{%
\begin{titlepage}
\thispagestyle{empty}
\vspace*{1.0in}
{\color{Navy}\rule{\linewidth}{2.2pt}}\\[0.55in]
{\Huge\bfseries\color{Navy} #1\par}
\vspace{0.18in}
{\Large\color{Teal} #2\par}
\vspace{0.5in}
\begin{tcolorbox}[colback=LightGray,colframe=Gray,boxrule=0.5pt,arc=2mm]
\textbf{Purpose.} #3
\end{tcolorbox}
\vfill
{\large Prepared from the revised edition of \textit{The Intelligent Investor}, with Jason Zweig's commentary.\par}
\vspace{0.12in}
{\small #4\par}
\vspace{0.25in}
{\small\color{Gray}\parbox{0.90\textwidth}{\centering Educational material only. Historical thresholds and examples are preserved for study; they are not individualized investment, tax, or legal advice.}}\par
\vspace{0.35in}
{\color{Navy}\rule{\linewidth}{2.2pt}}
\end{titlepage}
}

\newcommand{\DocShortTitle}{Cornell Notes}
\begin{document}
\makedoccover{The Intelligent Investor - Cornell-Style Chapter Notes}{Introduction, Chapters 1--20, Postscript, and Buffett appendix}{Provides chapter-by-chapter cues, notes, summaries, and self-test prompts.}{Standalone LaTeX study document}
\tableofcontents\newpage

\begin{cautionbox}[Source and modernization boundary]
The underlying source is the fourth revised edition, with Benjamin Graham's text updated in 1971--1972 and later commentary by Jason Zweig. Historical yields, tax rules, company examples, market levels, and dollar thresholds are therefore period-specific. This document preserves the source's framework and labels later implementation ideas separately rather than silently treating historical figures as current facts.
\end{cautionbox}

\section{How to Use These Cornell Notes}
Each chapter is organized into three components: \textbf{Cues / Questions} for recall, \textbf{Notes} for the substantive answer, and a \textbf{Summary} that compresses the chapter into one decision rule. The notes paraphrase the source; they are intended for study and retrieval, not as a substitute for reading the chapter and commentary.

\section{Introduction: What This Book Expects to Accomplish}
\begin{longtable}{@{}p{0.29\linewidth}p{0.66\linewidth}@{}}\toprule\textbf{Cue / Question} & \textbf{Cornell Notes}\\\midrule\endfirsthead\toprule\textbf{Cue / Question} & \textbf{Cornell Notes}\\\midrule\endhead
What is the book trying to teach? & The objective is an investment policy for nonprofessionals, emphasizing principles and attitudes more than exhaustive security-analysis technique. \\
Why study market history? & Historical market behavior is used to inoculate the investor against claims that current conditions are unprecedented or that easy riches are dependable. \\
What is the investor's main adversary? & The source repeatedly treats temperament as a first-order investment variable: enthusiasm, crowd behavior, and confidence in forecasts can be more dangerous than lack of intelligence. \\
Why ask 'How much?' & Price must remain connected to what is received in return. A good business can still be a poor investment at an excessive price. \\
\bottomrule\end{longtable}
\begin{principlebox}[Chapter summary]Build a policy that survives both market volatility and your own emotional reactions; do not confuse a plausible story with a satisfactory purchase price.\end{principlebox}
\subsection*{Self-test prompts}\begin{itemize}\item Explain the chapter's main failure mode without using the chapter title.\item State one quantitative or procedural rule from the chapter.\item Identify one way a modern investor could misapply the chapter.\end{itemize}

\section{Chapter 1 - Investment versus Speculation}
\begin{longtable}{@{}p{0.29\linewidth}p{0.66\linewidth}@{}}\toprule\textbf{Cue / Question} & \textbf{Cornell Notes}\\\midrule\endfirsthead\toprule\textbf{Cue / Question} & \textbf{Cornell Notes}\\\midrule\endhead
What separates investing from speculation? & Investment requires analysis, protection of principal, and an adequate return; operations failing those conditions are speculative. \\
What should a defensive investor expect? & The defensive investor aims for satisfactory results without constant effort, while the enterprising investor seeks additional return only through additional intelligent work. \\
What is the role of a speculative account? & If speculation is undertaken, it should be consciously limited and kept separate from the investment program rather than blended into it. \\
\bottomrule\end{longtable}
\begin{principlebox}[Chapter summary]Classify the operation before judging it. The label determines the required controls, sizing, and expectations.\end{principlebox}
\subsection*{Self-test prompts}\begin{itemize}\item Explain the chapter's main failure mode without using the chapter title.\item State one quantitative or procedural rule from the chapter.\item Identify one way a modern investor could misapply the chapter.\end{itemize}

\section{Chapter 2 - The Investor and Inflation}
\begin{longtable}{@{}p{0.29\linewidth}p{0.66\linewidth}@{}}\toprule\textbf{Cue / Question} & \textbf{Cornell Notes}\\\midrule\endfirsthead\toprule\textbf{Cue / Question} & \textbf{Cornell Notes}\\\midrule\endhead
Do stocks automatically protect against inflation? & The historical discussion rejects the idea that common stocks are an automatic or perfect hedge against inflation. \\
Why retain multiple asset types? & The portfolio compromise between equities and high-quality bonds reflects uncertainty about both inflation and business outcomes. \\
What did later commentary add? & The later commentary introduces inflation-protected securities as a direct tool for protecting part of fixed-income purchasing power. \\
\bottomrule\end{longtable}
\begin{principlebox}[Chapter summary]Inflation risk is real, but the response should be diversified and instrument-specific rather than a blanket assumption that stocks always solve it.\end{principlebox}
\subsection*{Self-test prompts}\begin{itemize}\item Explain the chapter's main failure mode without using the chapter title.\item State one quantitative or procedural rule from the chapter.\item Identify one way a modern investor could misapply the chapter.\end{itemize}

\section{Chapter 3 - A Century of Stock-Market History}
\begin{longtable}{@{}p{0.29\linewidth}p{0.66\linewidth}@{}}\toprule\textbf{Cue / Question} & \textbf{Cornell Notes}\\\midrule\endfirsthead\toprule\textbf{Cue / Question} & \textbf{Cornell Notes}\\\midrule\endhead
Why study long market series? & The chapter studies long-run price cycles and relates market prices to earnings and dividends. \\
Which relationships matter? & Valuation levels change dramatically across decades; therefore current multiples must be interpreted in historical and fundamental context. \\
What is the forecasting lesson? & The purpose of history is preparedness, not precise prediction. It shows the range of outcomes an investor must be financially and psychologically able to tolerate. \\
\bottomrule\end{longtable}
\begin{principlebox}[Chapter summary]History provides priors and guardrails; it does not remove uncertainty.\end{principlebox}
\subsection*{Self-test prompts}\begin{itemize}\item Explain the chapter's main failure mode without using the chapter title.\item State one quantitative or procedural rule from the chapter.\item Identify one way a modern investor could misapply the chapter.\end{itemize}

\section{Chapter 4 - General Portfolio Policy: The Defensive Investor}
\begin{longtable}{@{}p{0.29\linewidth}p{0.66\linewidth}@{}}\toprule\textbf{Cue / Question} & \textbf{Cornell Notes}\\\midrule\endfirsthead\toprule\textbf{Cue / Question} & \textbf{Cornell Notes}\\\midrule\endhead
How should stocks and bonds be combined? & A broad 25\%--75\% equity band with a 50/50 reference allocation is proposed for the defensive investor. \\
Why not maximize yield? & The desired return should reflect intelligent effort, not a simple willingness to take more risk. \\
What does risk mean here? & Portfolio construction is a compromise among purchasing-power risk, market risk, income needs, and behavioral stability. \\
\bottomrule\end{longtable}
\begin{principlebox}[Chapter summary]Choose an allocation you can maintain. A theoretically optimal portfolio that triggers bad behavior is not operationally optimal.\end{principlebox}
\subsection*{Self-test prompts}\begin{itemize}\item Explain the chapter's main failure mode without using the chapter title.\item State one quantitative or procedural rule from the chapter.\item Identify one way a modern investor could misapply the chapter.\end{itemize}

\section{Chapter 5 - The Defensive Investor and Common Stocks}
\begin{longtable}{@{}p{0.29\linewidth}p{0.66\linewidth}@{}}\toprule\textbf{Cue / Question} & \textbf{Cornell Notes}\\\midrule\endfirsthead\toprule\textbf{Cue / Question} & \textbf{Cornell Notes}\\\midrule\endhead
Why must a defensive investor own stocks? & Common stocks provide participation in business growth, but the purchase price remains critical. \\
What are the four basic rules? & The basic rules emphasize diversification, large/prominent and conservatively financed companies, long dividend records, and valuation limits relative to earnings. \\
How does dollar-cost averaging help? & Dollar-cost averaging converts market volatility into a systematic purchasing process and reduces the temptation to concentrate purchases at euphoric prices. \\
\bottomrule\end{longtable}
\begin{principlebox}[Chapter summary]Make the defensive process boring on purpose: diversified holdings, reasonable prices, automatic contributions, limited turnover.\end{principlebox}
\subsection*{Self-test prompts}\begin{itemize}\item Explain the chapter's main failure mode without using the chapter title.\item State one quantitative or procedural rule from the chapter.\item Identify one way a modern investor could misapply the chapter.\end{itemize}

\section{Chapter 6 - Enterprising Investor: Negative Approach}
\begin{longtable}{@{}p{0.29\linewidth}p{0.66\linewidth}@{}}\toprule\textbf{Cue / Question} & \textbf{Cornell Notes}\\\midrule\endfirsthead\toprule\textbf{Cue / Question} & \textbf{Cornell Notes}\\\midrule\endhead
What should the enterprising investor avoid? & The enterprising investor begins from the same high-quality base as the defensive investor. \\
Why are new issues dangerous? & The chapter warns against inferior bonds and preferreds purchased merely for yield, foreign-government issues without a compelling margin, and new issues sold into promotional conditions. \\
Why is extra yield not automatically extra value? & Complexity and a higher coupon do not compensate for weak credit or an inadequate purchase price. \\
\bottomrule\end{longtable}
\begin{principlebox}[Chapter summary]The first source of excess return is often what you refuse to buy.\end{principlebox}
\subsection*{Self-test prompts}\begin{itemize}\item Explain the chapter's main failure mode without using the chapter title.\item State one quantitative or procedural rule from the chapter.\item Identify one way a modern investor could misapply the chapter.\end{itemize}

\section{Chapter 7 - Enterprising Investor: Positive Side}
\begin{longtable}{@{}p{0.29\linewidth}p{0.66\linewidth}@{}}\toprule\textbf{Cue / Question} & \textbf{Cornell Notes}\\\midrule\endfirsthead\toprule\textbf{Cue / Question} & \textbf{Cornell Notes}\\\midrule\endhead
Where can active value come from? & One approach is to examine large companies temporarily out of favor, where scale and resources improve recovery odds. \\
Why can unpopular large companies be attractive? & A bargain issue is one whose analytically indicated value is materially above market price; Graham suggests at least a 50\% value premium over price as a concrete benchmark. \\
What is a bargain issue? & Net-current-asset opportunities represent an asset-based form of bargain where fixed assets are assigned little or no purchase value. \\
\bottomrule\end{longtable}
\begin{principlebox}[Chapter summary]Active investing requires a demonstrable price/value discrepancy, not merely an unpopular stock.\end{principlebox}
\subsection*{Self-test prompts}\begin{itemize}\item Explain the chapter's main failure mode without using the chapter title.\item State one quantitative or procedural rule from the chapter.\item Identify one way a modern investor could misapply the chapter.\end{itemize}

\section{Chapter 8 - The Investor and Market Fluctuations}
\begin{longtable}{@{}p{0.29\linewidth}p{0.66\linewidth}@{}}\toprule\textbf{Cue / Question} & \textbf{Cornell Notes}\\\midrule\endfirsthead\toprule\textbf{Cue / Question} & \textbf{Cornell Notes}\\\midrule\endhead
What is timing versus pricing? & Timing predicts market direction; pricing decides whether a security is attractive relative to value. \\
How should quotations be interpreted? & Market fluctuations are inevitable and can be useful, but they become dangerous when they pull the investor into speculative behavior. \\
What does Mr. Market teach? & Mr. Market personifies quotations as optional offers from an unstable counterparty, preserving the investor's independence of judgment. \\
\bottomrule\end{longtable}
\begin{principlebox}[Chapter summary]Use volatility to improve terms, not to outsource your judgment.\end{principlebox}
\subsection*{Self-test prompts}\begin{itemize}\item Explain the chapter's main failure mode without using the chapter title.\item State one quantitative or procedural rule from the chapter.\item Identify one way a modern investor could misapply the chapter.\end{itemize}

\section{Chapter 9 - Investing in Investment Funds}
\begin{longtable}{@{}p{0.29\linewidth}p{0.66\linewidth}@{}}\toprule\textbf{Cue / Question} & \textbf{Cornell Notes}\\\midrule\endfirsthead\toprule\textbf{Cue / Question} & \textbf{Cornell Notes}\\\midrule\endhead
Can superior funds be identified in advance? & The chapter evaluates open-end/closed-end, load/no-load, balanced/equity, and performance-oriented funds while questioning the persistence of exceptional manager performance. \\
What fund characteristics matter? & Fees, sales loads, diversification, and the risk of selecting recent winners matter because investor return is net of cost and behavior. \\
What does the commentary favor? & The later commentary presents low-cost index funds as a particularly strong defensive implementation. \\
\bottomrule\end{longtable}
\begin{principlebox}[Chapter summary]When manager selection is uncertain, broad diversification and low cost are controllable sources of advantage.\end{principlebox}
\subsection*{Self-test prompts}\begin{itemize}\item Explain the chapter's main failure mode without using the chapter title.\item State one quantitative or procedural rule from the chapter.\item Identify one way a modern investor could misapply the chapter.\end{itemize}

\section{Chapter 10 - The Investor and His Advisers}
\begin{longtable}{@{}p{0.29\linewidth}p{0.66\linewidth}@{}}\toprule\textbf{Cue / Question} & \textbf{Cornell Notes}\\\midrule\endfirsthead\toprule\textbf{Cue / Question} & \textbf{Cornell Notes}\\\midrule\endhead
What should an adviser reasonably provide? & Professional advice is most defensible when it protects against mistakes and delivers the normal results the invested capital is entitled to earn. \\
When can unconventional advice be accepted? & If the investor relies heavily on others, the recommended program should remain conservative unless the adviser is exceptionally well known and understood. \\
What duty stays with the investor? & Trustworthiness, competence, and fees cannot be delegated away; the investor retains responsibility for selecting the adviser. \\
\bottomrule\end{longtable}
\begin{principlebox}[Chapter summary]Hire process and judgment, not promises of extraordinary return.\end{principlebox}
\subsection*{Self-test prompts}\begin{itemize}\item Explain the chapter's main failure mode without using the chapter title.\item State one quantitative or procedural rule from the chapter.\item Identify one way a modern investor could misapply the chapter.\end{itemize}

\section{Chapter 11 - Security Analysis for the Lay Investor}
\begin{longtable}{@{}p{0.29\linewidth}p{0.66\linewidth}@{}}\toprule\textbf{Cue / Question} & \textbf{Cornell Notes}\\\midrule\endfirsthead\toprule\textbf{Cue / Question} & \textbf{Cornell Notes}\\\midrule\endhead
What is the key bond test? & For bonds and preferreds, fixed-charge coverage over several years and in a weak year is central to safety analysis. \\
Why be cautious with high-growth valuation? & Common-stock valuation depends on earning power, financial strength, dividends, and growth, but precision falls as forecasts extend further into the future. \\
What is the simplified growth formula? & The simplified historical formula $V=E(8.5+2g)$ is presented as an approximation, not a true-value machine. \\
\bottomrule\end{longtable}
\begin{principlebox}[Chapter summary]Prefer robust simple arithmetic to fragile precision built on optimistic long-range forecasts.\end{principlebox}
\subsection*{Self-test prompts}\begin{itemize}\item Explain the chapter's main failure mode without using the chapter title.\item State one quantitative or procedural rule from the chapter.\item Identify one way a modern investor could misapply the chapter.\end{itemize}

\section{Chapter 12 - Things to Consider About Per-Share Earnings}
\begin{longtable}{@{}p{0.29\linewidth}p{0.66\linewidth}@{}}\toprule\textbf{Cue / Question} & \textbf{Cornell Notes}\\\midrule\endfirsthead\toprule\textbf{Cue / Question} & \textbf{Cornell Notes}\\\midrule\endhead
Why not trust one year's EPS? & A single year's EPS can be distorted by cyclicality, special charges, accounting classification, or capital-structure effects. \\
What accounting traps matter? & The investor should read footnotes, understand dilution, and examine adjustments that exclude supposedly nonrecurring costs. \\
How should earnings be normalized? & Multi-year normalized earning power is more useful than a headline number for valuation. \\
\bottomrule\end{longtable}
\begin{principlebox}[Chapter summary]Earnings quality is part of valuation; a precise multiple applied to poor-quality EPS is false precision.\end{principlebox}
\subsection*{Self-test prompts}\begin{itemize}\item Explain the chapter's main failure mode without using the chapter title.\item State one quantitative or procedural rule from the chapter.\item Identify one way a modern investor could misapply the chapter.\end{itemize}

\section{Chapter 13 - A Comparison of Four Listed Companies}
\begin{longtable}{@{}p{0.29\linewidth}p{0.66\linewidth}@{}}\toprule\textbf{Cue / Question} & \textbf{Cornell Notes}\\\midrule\endfirsthead\toprule\textbf{Cue / Question} & \textbf{Cornell Notes}\\\midrule\endhead
What does side-by-side comparison reveal? & The chapter compares companies across profitability, growth, financial structure, valuation, and asset backing rather than judging each in isolation. \\
Why can a good company still be overpriced? & A high-quality growth company can require years of exceptional performance to justify an already elevated price. \\
What supports a protected purchase? & Companies bought near asset value with adequate earnings and balance-sheet strength offer a different form of protection. \\
\bottomrule\end{longtable}
\begin{principlebox}[Chapter summary]Comparative analysis exposes how much optimism is embedded in the price.\end{principlebox}
\subsection*{Self-test prompts}\begin{itemize}\item Explain the chapter's main failure mode without using the chapter title.\item State one quantitative or procedural rule from the chapter.\item Identify one way a modern investor could misapply the chapter.\end{itemize}

\section{Chapter 14 - Stock Selection for the Defensive Investor}
\begin{longtable}{@{}p{0.29\linewidth}p{0.66\linewidth}@{}}\toprule\textbf{Cue / Question} & \textbf{Cornell Notes}\\\midrule\endfirsthead\toprule\textbf{Cue / Question} & \textbf{Cornell Notes}\\\midrule\endhead
What are the seven tests? & The seven tests cover size, financial condition, earnings stability, dividend record, earnings growth, P/E, and price/book. \\
How do P/E and P/B interact? & The combined rule $P/E \times P/B \leq 22.5$ allows some tradeoff between earnings valuation and asset valuation. \\
What does the screen guarantee? & The endnotes stress that no formula guarantees superior market performance; the objective is to improve the odds of getting value for money. \\
\bottomrule\end{longtable}
\begin{principlebox}[Chapter summary]A defensive screen is a quality-and-price filter, not a short-term performance predictor.\end{principlebox}
\subsection*{Self-test prompts}\begin{itemize}\item Explain the chapter's main failure mode without using the chapter title.\item State one quantitative or procedural rule from the chapter.\item Identify one way a modern investor could misapply the chapter.\end{itemize}

\section{Chapter 15 - Stock Selection for the Enterprising Investor}
\begin{longtable}{@{}p{0.29\linewidth}p{0.66\linewidth}@{}}\toprule\textbf{Cue / Question} & \textbf{Cornell Notes}\\\midrule\endfirsthead\toprule\textbf{Cue / Question} & \textbf{Cornell Notes}\\\midrule\endhead
How do active screens differ from defensive ones? & Enterprising screens loosen some defensive requirements but demand enough compensating cheapness and analysis. \\
What is the NCAV method? & The classic net-current-asset method buys below stripped-down working-capital value; Graham-Newman historically used two-thirds or less in many cases. \\
Why diversify heavily? & Because individual bargain issues can fail, the historical method relied on broad diversification and process consistency. \\
\bottomrule\end{longtable}
\begin{principlebox}[Chapter summary]Deep value is a portfolio strategy, not a license to concentrate in statistically cheap weak businesses.\end{principlebox}
\subsection*{Self-test prompts}\begin{itemize}\item Explain the chapter's main failure mode without using the chapter title.\item State one quantitative or procedural rule from the chapter.\item Identify one way a modern investor could misapply the chapter.\end{itemize}

\section{Chapter 16 - Convertible Issues and Warrants}
\begin{longtable}{@{}p{0.29\linewidth}p{0.66\linewidth}@{}}\toprule\textbf{Cue / Question} & \textbf{Cornell Notes}\\\midrule\endfirsthead\toprule\textbf{Cue / Question} & \textbf{Cornell Notes}\\\midrule\endhead
What does a convertible combine? & A convertible combines a senior security with an option to participate in common-stock appreciation. \\
What is the hidden tradeoff? & The investor often pays for that option through a lower coupon, weaker seniority, or a higher price; therefore the package must be decomposed economically. \\
Why scrutinize warrants? & Large warrant structures can dilute common shareholders and signal a top-heavy or promotional capital structure. \\
\bottomrule\end{longtable}
\begin{principlebox}[Chapter summary]Do not let hybrid packaging hide the economics of yield, credit quality, conversion value, and dilution.\end{principlebox}
\subsection*{Self-test prompts}\begin{itemize}\item Explain the chapter's main failure mode without using the chapter title.\item State one quantitative or procedural rule from the chapter.\item Identify one way a modern investor could misapply the chapter.\end{itemize}

\section{Chapter 17 - Four Extremely Instructive Case Histories}
\begin{longtable}{@{}p{0.29\linewidth}p{0.66\linewidth}@{}}\toprule\textbf{Cue / Question} & \textbf{Cornell Notes}\\\midrule\endfirsthead\toprule\textbf{Cue / Question} & \textbf{Cornell Notes}\\\midrule\endhead
What failed at Penn Central? & Penn Central illustrates ignored coverage and operating warnings. \\
What failed at LTV? & Ling-Temco-Vought illustrates acquisition-driven expansion funded by rapidly rising debt. \\
What do NVF and AAA add? & NVF highlights leverage and accounting complexity in an outsized acquisition; AAA illustrates speculative financing of a small fashionable concept that lacked durable economics. \\
\bottomrule\end{longtable}
\begin{principlebox}[Chapter summary]Perform elementary credit, balance-sheet, and earnings tests before accepting a growth or acquisition narrative.\end{principlebox}
\subsection*{Self-test prompts}\begin{itemize}\item Explain the chapter's main failure mode without using the chapter title.\item State one quantitative or procedural rule from the chapter.\item Identify one way a modern investor could misapply the chapter.\end{itemize}

\section{Chapter 18 - A Comparison of Eight Pairs of Companies}
\begin{longtable}{@{}p{0.29\linewidth}p{0.66\linewidth}@{}}\toprule\textbf{Cue / Question} & \textbf{Cornell Notes}\\\midrule\endfirsthead\toprule\textbf{Cue / Question} & \textbf{Cornell Notes}\\\midrule\endhead
Why compare pairs? & Paired analysis isolates differences in financial structure, operating history, valuation, and corporate policy among superficially similar securities. \\
What patterns recur? & Rapid growth, conglomeration, financial engineering, or popularity can create prices that outrun business protection. \\
What does competition do to high returns? & The endnotes emphasize a competitive-economics lesson: unusually high returns attract competitors, which can erode those returns. \\
\bottomrule\end{longtable}
\begin{principlebox}[Chapter summary]Compare alternatives, not just stories. Opportunity cost is visible when two similar-looking businesses have radically different financial quality and price.\end{principlebox}
\subsection*{Self-test prompts}\begin{itemize}\item Explain the chapter's main failure mode without using the chapter title.\item State one quantitative or procedural rule from the chapter.\item Identify one way a modern investor could misapply the chapter.\end{itemize}

\section{Chapter 19 - Shareholders and Managements: Dividend Policy}
\begin{longtable}{@{}p{0.29\linewidth}p{0.66\linewidth}@{}}\toprule\textbf{Cue / Question} & \textbf{Cornell Notes}\\\midrule\endfirsthead\toprule\textbf{Cue / Question} & \textbf{Cornell Notes}\\\midrule\endhead
What should retained earnings accomplish? & Retained earnings are justified when management can deploy them productively for shareholders. \\
When are dividends preferable? & Where reinvestment is weak, distributing cash can be superior to retaining it merely to enlarge the enterprise. \\
How should repurchases be judged? & The commentary extends the ownership lens to executive compensation and repurchases, favoring buybacks when shares are cheap rather than when management options make repurchases convenient. \\
\bottomrule\end{longtable}
\begin{principlebox}[Chapter summary]Treat capital allocation as evidence of management quality.\end{principlebox}
\subsection*{Self-test prompts}\begin{itemize}\item Explain the chapter's main failure mode without using the chapter title.\item State one quantitative or procedural rule from the chapter.\item Identify one way a modern investor could misapply the chapter.\end{itemize}

\section{Chapter 20 - Margin of Safety}
\begin{longtable}{@{}p{0.29\linewidth}p{0.66\linewidth}@{}}\toprule\textbf{Cue / Question} & \textbf{Cornell Notes}\\\midrule\endfirsthead\toprule\textbf{Cue / Question} & \textbf{Cornell Notes}\\\midrule\endhead
What is the central concept? & Margin of safety is the buffer between what must go right and what can go wrong. \\
How does it differ for bonds and stocks? & For bonds it appears as excess coverage of fixed charges; for stocks it can appear as earning power and value sufficiently above the purchase price or alternative yield. \\
Why does diversification still matter? & Diversification is complementary: even sound analysis can be wrong, so portfolio construction prevents one error from becoming fatal. \\
\bottomrule\end{longtable}
\begin{principlebox}[Chapter summary]Require a cushion against analytical error; do not make success depend on a single forecast being exactly right.\end{principlebox}
\subsection*{Self-test prompts}\begin{itemize}\item Explain the chapter's main failure mode without using the chapter title.\item State one quantitative or procedural rule from the chapter.\item Identify one way a modern investor could misapply the chapter.\end{itemize}

\section{Postscript - Preparation, Courage, and Uncertainty}
\begin{longtable}{@{}p{0.29\linewidth}p{0.66\linewidth}@{}}\toprule\textbf{Cue / Question} & \textbf{Cornell Notes}\\\midrule\endfirsthead\toprule\textbf{Cue / Question} & \textbf{Cornell Notes}\\\midrule\endhead
What does the GEICO episode teach? & Graham describes one unusually successful investment that contributed more than many years of routine decisions. \\
Was concentration the general rule? & The commentary explicitly warns that such a concentrated success should not be generalized; ordinary investors need diversification because they cannot rely on finding a similarly exceptional opportunity. \\
What is the role of uncertainty? & The broader lesson is that preparation creates the ability to recognize an unusual opportunity, while courage is required to act when uncertainty remains. \\
\bottomrule\end{longtable}
\begin{principlebox}[Chapter summary]Prepare broadly, act selectively, and do not confuse one extraordinary case with a default portfolio rule.\end{principlebox}
\subsection*{Self-test prompts}\begin{itemize}\item Explain the chapter's main failure mode without using the chapter title.\item State one quantitative or procedural rule from the chapter.\item Identify one way a modern investor could misapply the chapter.\end{itemize}

\section{Appendix - The Superinvestors of Graham-and-Doddsville}
\begin{longtable}{@{}p{0.29\linewidth}p{0.66\linewidth}@{}}\toprule\textbf{Cue / Question} & \textbf{Cornell Notes}\\\midrule\endfirsthead\toprule\textbf{Cue / Question} & \textbf{Cornell Notes}\\\midrule\endhead
What question does Buffett test? & Buffett challenges the idea that repeated value-investing outperformance can be dismissed automatically as random luck. \\
What is common among the investors? & The investors discussed differ in holdings and style but share a price-versus-value discipline and an emphasis on margin of safety. \\
What does this add to the framework? & The appendix reinforces process consistency and independent reasoning rather than imitation of any single portfolio. \\
\bottomrule\end{longtable}
\begin{principlebox}[Chapter summary]The common factor is a disciplined valuation philosophy, not a shared list of stocks.\end{principlebox}
\subsection*{Self-test prompts}\begin{itemize}\item Explain the chapter's main failure mode without using the chapter title.\item State one quantitative or procedural rule from the chapter.\item Identify one way a modern investor could misapply the chapter.\end{itemize}

\section*{Primary Source}
Benjamin Graham, \textit{The Intelligent Investor: A Book of Practical Counsel}, revised edition, updated with commentary by Jason Zweig, preface and appendix material by Warren E. Buffett. Chapter references in this document point to the structure of that edition.

\end{document}
